\documentclass{article}
\usepackage[T1]{fontenc}
\usepackage{lmodern}
\usepackage{graphicx}
\usepackage{amsmath,amssymb}
\usepackage[a4paper,margin=2cm]{geometry}
\usepackage[absolute,overlay]{textpos}
\setlength{\TPHorizModule}{1mm}
\setlength{\TPVertModule}{1mm}
\usepackage[indonesian]{babel}
\usepackage{hyperref}
\setlength{\emergencystretch}{2em}
% unit: Halaman 30

\begin{document}

% === Halaman 30 (python/native) ===

30

4.  Cipherteks (ciphertext): pesan yang telah disandikan sehingga tidak 

bermakna lagi.

      Tujuan: agar pesan tidak dapat dibaca oleh pihak yang tidak berhak.

      Nama lain: kriptogram (cryptogram)

• Contoh:

 Plainteks:  culik anak itu jam 11 siang

 Cipherteks: t\textasciicircum{}\$gfUi89rewoFpfdWqLMp[uTcxZ

 Kriptogram: t\textasciicircum{}\$gfU  i89rewo  FpfdWqLM  p[uTcxZ

\end{document}
